% !TEX program = pdflatex
% !TEX root = paper2_GRSL_letter_supplement.tex
%% =====================================================================
%% SUPPLEMENT to the GRSL letter
%% "Sensor-Family Composition of the In-Situ Archive Changes the
%%  Apparent Bias and RMSD of Satellite Soil-Moisture Products"
%%
%% Supplementary analyses for the 5-page main letter:
%%   Table S1  full family composition + per-reference residual CIs
%%   Fig. S1   per-archive dielectric-free CRNS attribution
%%   Table S2  near-surface paired-sensor test
%%   Table S3/Fig. S2 real 2021 SMAP temporal validation check
%%   Sec. S4   temporal-window invariance (ERA5 2020-23 vs 2023-only)
%%   Sec. S5   soil-texture control (clay/bulk-density mixed model)
%% =====================================================================
\documentclass[11pt]{article}
\usepackage[utf8]{inputenc}
\usepackage{amsmath,amssymb}
\usepackage{graphicx}
\usepackage{booktabs}
\usepackage[margin=2.4cm]{geometry}
\usepackage{url}
\setlength{\emergencystretch}{3em}

\graphicspath{%
  {paper2_GRSL_code_package_20260708/figures/main/}%
  {paper2_GRSL_code_package_20260708/figures/supplement/}%
  {figures/main/}{figures/supplement/}%
  {../figures/main/}{../figures/supplement/}}

\renewcommand{\thetable}{S\arabic{table}}
\renewcommand{\thefigure}{S\arabic{figure}}
\renewcommand{\thesection}{S\arabic{section}}

\title{Supplementary Material\\
\large Sensor-Family Composition of the In-Situ Archive Changes the
Apparent Bias and RMSD of Satellite Soil-Moisture Products}
\author{Chang-Hwan Park\\
\small Department of Civil Systems Engineering, Ajou University\\
\small Suwon, Republic of Korea\\
\small Department of Civil Systems Engineering, Ajou University in Tashkent\\
\small Tashkent, Uzbekistan}
\date{July 8, 2026}

\begin{document}
\maketitle

\section{Archive composition and per-reference residuals}
Table~\ref{tab:S1} gives the full manufacturer-family composition of the
retained 2020–2023 ISMN archive and the mean ISMN$-$reference residual
for each family against all three reference products, with 95\,\%
bootstrap confidence intervals. The HydraProbe-family ERA5-Land CI does
not overlap the Decagon/METER TDR CI, and the family swing
($0.17$~m$^3$\,m$^{-3}$) is an order of magnitude larger than the
within-family disagreement among references.

\begin{table}[h!]
\centering
\caption{ISMN 2020–2023 composition and mean ISMN$-$reference residual
(m$^3$\,m$^{-3}$) by manufacturer family, 95\,\% bootstrap CI (5{,}000
resamples). Residual means use the valid-coverage matchup set; the
composition count and the ERA5 residual count differ because not every
station-sensor has valid coverage in every product (e.g., 1{,}231 of
1{,}274 HydraProbe-family sensors have valid ERA5-Land coverage).}
\label{tab:S1}
\small
\resizebox{\textwidth}{!}{%
\begin{tabular}{lccccc}
\toprule
Sensor family & $n$ & Share (\%) & ERA5-Land & SMAP~L3\_E & SMAP-SMI\\
\midrule
HydraProbe family & 1{,}274 & 78.4 & $-0.097$ $[-0.103,-0.091]$ & $-0.025$ $[-0.035,-0.016]$ & $-0.037$ $[-0.043,-0.032]$\\
Decagon/METER TDR &   95 &  5.8 & $+0.068$ $[+0.028,+0.109]$ & $+0.006$ $[-0.041,+0.057]$ & $-0.074$ $[-0.103,-0.042]$\\
Campbell CS6xx    &   65 &  4.0 & $-0.065$ $[-0.088,-0.042]$ & $+0.019$ $[-0.001,+0.039]$ & $-0.024$ $[-0.050,+0.005]$\\
ThetaProbe        &   59 &  3.6 & $-0.039$ $[-0.061,-0.017]$ & $+0.001$ $[-0.025,+0.027]$ & $-0.001$ $[-0.028,+0.025]$\\
TRIME             &    2 &  0.1 & — & — & —\\
other             &  130 &  8.0 & $-0.056$ $[-0.079,-0.031]$ & $+0.003$ $[-0.026,+0.033]$ & $-0.022$ $[-0.046,+0.002]$\\
\midrule
Total             & 1{,}625 & 100.0 & & & \\
\bottomrule
\end{tabular}
}
\end{table}

\section{Per-archive dielectric-free CRNS attribution}
The pooled cross-check reported in the main text (92
CRNS$\leftrightarrow$ISMN pairs, Welch's $t=+4.07$,
$p=2.3\times10^{-4}$) decomposes by archive as follows.
\emph{COSMOS-USA} (reprocessed CRNS$-$ERA5-Land bias table of Zheng
et al.~\cite{ZhengHESS2024}, matched within $0.2^\circ$, $n=56$):
HydraProbe
$+0.142$~m$^3$\,m$^{-3}$ $[+0.114,+0.170]$ vs non-HydraProbe $-0.067$
$[-0.174,+0.037]$; Welch's $t=+3.61$, $p=1.4\times10^{-3}$.
\emph{COSMOS-Europe}~\cite{BogenaCOSMOSEurope2022} (matched within
$0.1^\circ$, $n=31$): HydraProbe
$+0.058$ $[+0.038,+0.078]$ vs non-HydraProbe $+0.010$ $[-0.016,+0.040]$;
$t=+2.60$, $p=0.019$. \emph{COSMOS-UK}~\cite{Cooper2021COSMOSUK}
on-site CRNS$-$TDT (same-footprint dielectric probe, $n=49$):
$+0.035$ $[+0.010,+0.061]$, excluded from the pooled within-archive
family contrast because its design has no non-HydraProbe channel.
\emph{CosmOz Australia}~\cite{Hawdon2014CosmOz} ($n=5$):
CRNS$-$HydraProbe $+0.162$ ($n=2$), CRNS$-$Campbell CS6xx $-0.043$
($n=3$)—directionally consistent but too small for an independent
test. Leave-one-archive deletion: dropping USA gives $t=+3.91$,
$p=6.3\times10^{-4}$; dropping Europe $t=+4.05$, $p=3.7\times10^{-4}$;
dropping Australia $t=+3.68$, $p=7.8\times10^{-4}$. Figure~\ref{fig:S1}
summarizes the archive-level contrasts.

The COSMOS-USA row is a synthesized residual: the Zheng et al.
CRNS$-$ERA5-Land bias is differenced against the station-matched
ISMN$-$ERA5-Land residual from this work. The two ERA5-Land overlays are
not redistributed as common-cell grids in this compact review package,
so the USA row should be interpreted as a supporting rather than
sole-identifying contrast. The USA-excluded leave-one-archive result
above is therefore the conservative check used in the main text.

\begin{figure}[h!]
\centering
\includegraphics[width=0.9\textwidth]{phase2_combined_attribution.png}
\caption{Per-archive dielectric-free CRNS attribution. Each row is the
mean (CRNS $-$ dielectric probe) residual for one archive-family
comparison; error bars are 95\,\% bootstrap CI. The neutron reference
is systematically wetter than the HydraProbe family and closer to
non-HydraProbe families across three continents; Australia is shown for
direction only.}
\label{fig:S1}
\end{figure}

\section{Near-surface constraint for the paired-sensor test}
The main-text paired test uses $n=21$ stations that host both a
HydraProbe-family and a non-HydraProbe sensor in the same ISMN grid
cell; both members are in the retained near-surface layer
(depth-from $\le0.15$~m). Table~\ref{tab:S2} reports the paired
family contrast used in the main text. The inference is therefore
restricted to co-located near-surface pairs.

\begin{table}[h!]
\centering
\caption{Paired-sensor offset (non-HydraProbe $-$ HydraProbe) for
co-located near-surface pairs.}
\label{tab:S2}
\small
\resizebox{0.82\textwidth}{!}{%
\begin{tabular}{lcccc}
\toprule
Pair subset & pairs & mean diff (m$^3$m$^{-3}$) & $t$ & $p$\\
\midrule
All co-located pairs & 21 & $+0.075$ $[+0.031,+0.119]$ & $+3.19$ & $4.7\times10^{-3}$\\
\bottomrule
\end{tabular}
}
\end{table}

\section{Real 2021 SMAP temporal validation check}
The main letter focuses on mean-state bias and RMSD. To connect the
family-composition result to standard satellite calibration/validation
language, we also computed a real overpass-matched daily matchup for
2021 using raw ISMN \texttt{.stm} hourly files and SMAP L3 SM\_P
descending/AM granules. The filters were ISMN quality flag "G",
depth-from $\le0.15$~m, recommended-quality SMAP retrievals, and at
least 20 matched days per station-sensor. The resulting table contains
785 station-sensors and 78{,}600 matched station-days.

\begin{table}[h!]
\centering
\caption{Real 2021 overpass-matched temporal metrics by manufacturer
family against SMAP L3 SM\_P (AM). Bias is ISMN minus SMAP; values are
station-sensor medians; brackets give 95\,\% network-cluster bootstrap
CI. ThetaProbe ($n=3$) is shown for completeness only.}
\label{tab:S3}
\small
\resizebox{\textwidth}{!}{%
\begin{tabular}{lccccc}
\toprule
Family & Stations & Bias & ubRMSE & $R$ & Anomaly $R$\\
\midrule
HydraProbe & 637 & $-0.032$ $[-0.099,-0.017]$ & $0.062$ $[0.051,0.074]$ & $0.317$ $[0.150,0.443]$ & $0.329$ $[0.088,0.465]$\\
Campbell CS6xx & 58 & $-0.107$ $[-0.151,+0.099]$ & $0.058$ $[0.032,0.064]$ & $0.166$ $[0.087,0.622]$ & $0.140$ $[-0.036,0.657]$\\
Decagon/METER TDR & 57 & $+0.058$ $[+0.031,+0.139]$ & $0.048$ $[0.028,0.068]$ & $0.622$ $[0.441,0.647]$ & $0.423$ $[0.264,0.635]$\\
other & 30 & $-0.071$ $[-0.085,-0.033]$ & $0.078$ $[0.056,0.080]$ & $0.055$ $[0.001,0.436]$ & $0.059$ $[-0.016,0.283]$\\
ThetaProbe & 3 & $+0.195$ & $0.030$ & $0.639$ & $0.608$\\
\bottomrule
\end{tabular}
}
\end{table}

\begin{figure}[h!]
\centering
\includegraphics[width=\textwidth]{family_temporal_real2021.png}
\caption{Real 2021 overpass-matched SMAP L3 SM\_P (AM) temporal
validation by sensor family. Bars are family medians and error bars are
network-cluster 95\,\% CI.}
\label{fig:S2}
\end{figure}

The HydraProbe-versus-non-HydraProbe contrast is
$-0.032$ versus $-0.036$~m$^3$\,m$^{-3}$ for bias
($p=9.8\times10^{-3}$), $0.062$ versus
$0.057$~m$^3$\,m$^{-3}$ for ubRMSE ($p=8.2\times10^{-4}$),
$0.317$ versus $0.273$ for $R$ ($p=0.70$), and $0.329$ versus $0.179$
for anomaly $R$ ($p=0.055$). Because this check uses the 36-km SMAP
SM\_P product, one year, AM overpasses only, and point-to-grid matchups,
it is interpreted as a standard-metric diagnostic rather than as a
replacement for the mean-state SPL3SMP\_E result in the main text.
The pooled non-HydraProbe bias masks family heterogeneity:
Decagon/METER TDR is wetter than HydraProbe ($+0.058$ versus
$-0.032$~m$^3$\,m$^{-3}$), whereas Campbell CS6xx and other records are
drier ($-0.107$ and $-0.071$~m$^3$\,m$^{-3}$).

\section{Temporal-window invariance}
Because SMAP~L3\_E was extracted for 2023 only while ERA5-Land and
SMAP-SMI span 2020–2023, we recomputed the ERA5-Land swvl1 (0–7~cm)
grand mean for (i)~the full 2020–2023 window and (ii)~2023 only, using
an identical $0.1^\circ$ nearest-neighbour protocol. The HydraProbe-family
mean residual changes by $0.004$~m$^3$\,m$^{-3}$ ($3.9\,\%$) between the
two windows; all five families change by $\le0.012$~m$^3$\,m$^{-3}$; and
the driest$\rightarrow$wettest family ranking (HydraProbe $\rightarrow$
Decagon/METER $\rightarrow$ other $\rightarrow$ ThetaProbe $\rightarrow$
Campbell~CS6xx) is identical in both windows. The reported family
ordering is therefore not an artifact of the ERA5 2020–2023 versus
SMAP~L3\_E 2023-only window choice.

\section{Soil-texture control}
Soil organic carbon (SOC; g\,kg$^{-1}$) was log-transformed before
regression. A single ordinary-least-squares regression of the pooled
three-reference mean residual on $\log_{10}$(SOC) gives a slope of
$-0.078$~m$^3$\,m$^{-3}$ per decade ($p=1.7\times10^{-8}$, $n=786$).
At face value this reads as an organic-matter effect, but it is not
robust to texture control. In a linear mixed-effects model with a
random intercept for ISMN network and fixed-effect covariates for
SoilGrids~2.0~\cite{Poggio2021} clay fraction and bulk density, sensor
depth-from, and $|\text{latitude}|$, implemented with
Statsmodels~\cite{Seabold2010}, the $\log_{10}$(SOC) coefficient
collapses in all three references (ERA5-Land $-0.029$, $p=0.11$;
SMAP~L3\_E $-0.009$, $p=0.77$; SMAP-SMI $+0.030$, $p=0.08$), while
the clay coefficient is highly significant in every reference
($p<10^{-9}$). A leave-one-network jackknife of the OLS slope stays
negative across all networks ($-0.074\pm0.011$; range $-0.096$ to
$-0.034$). The apparent organic-matter dependence is thus confounded
with clay fraction and with HydraProbe deployment geography, consistent
with a
sensor-family/texture confounder rather than a direct organic-matter
effect; the sensor-family ordering itself is unchanged by these
controls.

\begin{thebibliography}{9}
\bibitem{ZhengHESS2024} Y.~Zheng \emph{et al.}, "Evaluation of
reanalysis soil moisture products using cosmic ray neutron sensor
observations across the globe," \emph{Hydrol. Earth Syst. Sci.},
vol.~28, pp.~1999–2022, 2024, doi:10.5194/hess-28-1999-2024.
\bibitem{BogenaCOSMOSEurope2022} H.~R.~Bogena \emph{et al.},
"COSMOS-Europe: a European network of cosmic-ray neutron soil moisture
sensors," \emph{Earth Syst. Sci. Data}, vol.~14, pp.~1125–1151, 2022,
doi:10.5194/essd-14-1125-2022.
\bibitem{Cooper2021COSMOSUK} H.~M.~Cooper \emph{et al.}, "Soil
moisture estimates from the UK COSMOS-UK network," \emph{Earth Syst.
Sci. Data}, vol.~13, pp.~1737–1757, 2021,
doi:10.5194/essd-13-1737-2021.
\bibitem{Hawdon2014CosmOz} A.~Hawdon, D.~McJannet, and J.~Wallace,
"Calibration and correction procedures for cosmic-ray neutron soil
moisture probes across Australia," \emph{Water Resour. Res.}, vol.~50,
pp.~5029–5043, 2014, doi:10.1002/2013WR015138.
\bibitem{Poggio2021} L.~Poggio \emph{et al.}, "SoilGrids 2.0:
producing soil information for the globe with quantified spatial
uncertainty," \emph{SOIL}, vol.~7, pp.~217–240, 2021,
doi:10.5194/soil-7-217-2021.
\bibitem{Seabold2010} S.~Seabold and J.~Perktold, "Statsmodels:
econometric and statistical modeling with Python," in \emph{Proc. 9th
Python in Science Conf.}, 2010.
\end{thebibliography}

\end{document}
